\subsection*{A failure of the proposed graded recognition}

\statement{B.16}{Refuted}{The graded assertion of P7 fails in weight four}
Over \(\mathbb Q\), the common kernel of the four forgetting maps
from outer special derivations of the four-punctured sphere to
outer special derivations of three-punctured spheres is strictly
larger than the pure braid Lie algebra. In weight four its
dimension is six, whereas the braid subspace has dimension three.
An explicit element of the difference is given below.
\begin{proof}
Let \(L=L(x,y,z)\) be the free graded Lie algebra on three generators
in degree one and put \(x_4=-x-y-z\).
A tangential derivation of degree \(k>1\) is specified by
\[
 D(x_i)=[x_i,a_i],\qquad a_i\in L_k,\qquad 1\leq i\leq4,
\]
subject to \(\sum_iD(x_i)=0\). Subtracting the inner derivation
with common tangential conjugator \(a_4\) sets \(a_4=0\).
This representative is unique: the centralizer of the nonzero
degree-one element \(x_4\) in a free Lie algebra has no component
in degree \(k>1\). Thus outer special derivations identify with
\[
 \ker\left(L_k^{\,3}\longrightarrow L_{k+1},\
             (a_1,a_2,a_3)\longmapsto
             [x,a_1]+[y,a_2]+[z,a_3]\right).
\]
The arrow is surjective. Indeed the Jacobi identity rewrites
every Lie monomial of degree \(k+1\) as a sum of brackets with a
degree-one first entry. Witt's dimensions give
\(\dim L_4(3)=18\), \(\dim L_5(3)=48\), so the kernel has
dimension \(3\cdot18-48=6\).

For a three-punctured sphere the identical argument uses
\(L(x,y)\) and gives dimension
\(2\dim L_4(2)-\dim L_5(2)=2\cdot3-6=0\).
Therefore every one of the four maps has zero target in this
weight. The braid Lie algebra of \(P\Gamma_{0,4}\) is the free
Lie algebra on two degree-one generators, so its weight-four
dimension is three. To give an explicit witness, define
\[
\begin{split}
 a_1&=[x,[x,[y,z]]]-3[[x,y],[x,z]],\\
 a_2&=-[x,[x,[x,z]]],\\
 a_3&=[x,[x,[x,y]]],\qquad a_4=0.
\end{split}
\]
Then \([x,a_1]+[y,a_2]+[z,a_3]=0\), by Jacobi, so this is
a special derivation. Forgetting \(x\), \(y\), or \(z\) gives
the zero derivation; forgetting \(x_4\) means substituting
\(z=-x-y\) and gives an inner derivation with common
tangential conjugator \([x,[x,[x,y]]]\).

For the nonmembership check let
\[
\begin{aligned}
 A(x)&=[x,y],& A(y)&=[y,x],& A(z)&=0,\\
 B(x)&=0,& B(y)&=[y,z],& B(z)&=[z,y].
\end{aligned}
\]
These are the two infinitesimal pure braid generators.
The three weight-four Lyndon brackets in \(A,B\) have rank three
in the free associative-algebra model. Adjoining the displayed
derivation raises the rank to four. The program
\texttt{code/check\_special\_derivations.py} checks all these
identities and ranks over \(\mathbb Q\), using explicit Lyndon
bases and associative expansions. It also verifies all four
forgetting maps, including the inner conjugator for the fourth.
\end{proof}

\statement{B.17}{Refuted}{The full group-theoretic equality in P7 is false}
The common kernel of the three corrected resolvent maps strictly
contains \(\mathcal K\). The full counterexample is the nonempty
lift set of B.57, whose members induce the identity on all three
full filled targets but have a nontrivial tripod action detected
by the outer transposition of \(A_9\).
\begin{proof}
Apply B.57, proved below by an exact retained quotient and its
actual central refinement. This proof is separate from B.16.
Passing from a rational derivation to a full profinite
automorphism is not a map supplied by that computation; no
integration of B.16 is asserted or used.
\end{proof}
